% Generated draft: requires mathematical and source review.
\documentclass[11pt]{article}
\usepackage[margin=25mm]{geometry}
\usepackage{fontspec}
\setmainfont{Latin Modern Roman}
\usepackage{amsmath,amssymb,mathtools,unicode-math}
\setmathfont{Latin Modern Math}
\usepackage{enumitem,needspace,xurl,hyperref}
\setlength{\parindent}{0pt}
\setlength{\parskip}{5pt}
\setlength{\emergencystretch}{3em}
\newcommand{\N}{\mathbb{N}}
\newcommand{\Z}{\mathbb{Z}}
\newcommand{\Q}{\mathbb{Q}}
\newcommand{\R}{\mathbb{R}}
\newcommand{\C}{\mathbb{C}}
\newcommand{\F}{\mathbb{F}}
\newcommand{\A}{\mathbb{A}}
\newcommand{\PP}{\mathbb{P}}
\newcommand{\E}{\mathbb{E}}
\newcommand{\eps}{\varepsilon}
\newcommand{\dd}{\,\mathrm{d}}
\begin{document}
\section*{UnsolvedMath ID: 20002677}
\textbf{Source problem code:} catalogue entry
\subsection{Definitions and mathematical statement}
Let \((\mathcal M,\tau)\) be a finite tracial von Neumann algebra, and let \(x,y\) be free unit-variance circular elements, free from the \(\mathrm{II}_1\) algebra containing bounded \(S,T\). For an affiliated measurable operator use \(\|A\|_p=(\tau|A|^p)^{1/p}\), a quasi-norm when \(0<p<1\). Inverses in the source are affiliated inverses, not necessarily bounded resolvents.

The source asks whether \(x\) can replace \(xy^{-1}\) in its Haagerup--Schultz inverse-difference estimate for \(0<p<2/3\), with a constant depending only on \(p\). Its extraction obscures the norm subscripts and a displayed reference label. Reading the two printed \(p\)'s as norm subscripts gives the target
\[
 \|(S+x)^{-1}-(T+x)^{-1}\|_p\le C_p\|S-T\|_p.
\]
If the original right side is the operator norm instead, that is a distinct operator-valued question requiring primary-source confirmation. The scalar partial results below use \(S=sI,T=tI\), where both norms equal \(|s-t|\), so they do not resolve that ambiguity. An exponent on \(|s-t|\) is not silently substituted for the aid's Lipschitz reading.

The aid's scalar theorem is: for a unit-variance circular element \(c\), every \(0<q<1\) satisfies
\[
 \sup_{\lambda\in\C}\|(\lambda-c)^{-1}\|_q^q=M_q,
 \qquad
 M_q=\frac{2^{1-q}}\pi B\left(\frac{1-q}{2},\frac32\right),
\]
where \(B(a,b)=\int_0^1t^{a-1}(1-t)^{b-1}\,\dd t\) is the beta function. Equality in the supremum occurs at \(\lambda=0\). Consequently, for every \(0<p<1/2\) and \(s,t\in\C\),
\[
 \|(s+c)^{-1}-(t+c)^{-1}\|_p
 \le M_{2p}^{1/p}|s-t|.
\]
The constant comes from the resolvent identity and noncommutative H\"older with two \(L^{2p}\) inverse factors; it does not cover \(p\ge1/2\), since the required inverse moment is then outside \(q<1\).

The aid's sharp necessary threshold is separate. If a uniform scalar Lipschitz estimate at a given \(p>0\) holds for all \(s,t\), scalar difference quotients at zero and lower semicontinuity force \(c^{-2}\in L^p\). The supplied Fuss--Catalan hard-edge calculation gives
\[
 c^{-2}\in L^p\quad\Longleftrightarrow\quad p<2/3.
\]
Here \(c^{-2}=(c^2)^{-1}\) is not replaced by \(|c|^{-2}\), since \(c\) is nonnormal. Thus no such uniform scalar Lipschitz estimate is possible for \(p\ge2/3\). This shows necessity of the source's upper endpoint, not sufficiency in the remaining strip.

The interval \(1/2\le p<2/3\) even for scalar perturbations, and the general operator-valued \(S,T\) assertion in the source's intended norm convention, remain unresolved in the supplied aid. The scalar result does not settle the proposed circular replacement of the original free Cauchy-type perturbation.
\subsection{Short English statement}
The aid proves the circular inverse-difference estimate for scalar perturbations when \(p<1/2\), with an explicit inverse-moment constant. Scalar difference quotients exclude \(p\ge2/3\). The intermediate scalar strip and the full operator-valued question remain open within the supplied result, with a damaged right-norm convention still requiring review.
\subsection{Sources}
\begin{itemize}
\item Local registry, permanent UnsolvedMath ID \(20002677\), exact public code \texttt{catalogue entry}; canonical location \texttt{aim-probability-notes.json}, \texttt{notes[118]}.
\item AIM, \emph{Free Analysis}, problem list dated 24 August 2006, section 0.6, ``Invariant Subspaces for an Operator,'' question attributed to Haagerup; supplied reference \url{https://aimath.org/WWN/freeanalysis/freeanalysis.pdf}. The verification field reports prior visual checking of page 4 of the seven-page PDF to recover inverse placement and \(2/3\). No PDF was accessed here.
\item Haagerup--Schultz estimate, named in the source with an unresolved displayed label. The aid invokes a Fuss--Catalan hard-edge calculation but supplies no complete bibliographic references for either ingredient.
\item Supplied research attempt 1: \texttt{partial\_result}, run status \texttt{unknown}, rigor label \texttt{proved}, low novelty confidence; registry status \texttt{partially\_solved}, classification \texttt{PARTIAL-PROGRESS}, created and updated 14 August 2026. The aid is machine generated, not peer reviewed, and requires independent verification before citation.
\end{itemize}
\end{document}
